%%%%%%%%%%  scan idx 212  |  folio 199  %%%%%%%%%%
\section*{Appendice II. Définition de l'indice analytique d'un complexe elliptique relatif}

On se propose, dans cet appendice, de montrer que la notion de complexe parfait conduit à une définition simple et naturelle de «~l'indice analytique~» d'une famille continue d'opérateurs différentiels elliptiques, ou plus généralement d'un «~complexe elliptique relatif~». Pour le lecteur initié, disons que la définition donnée ici présente l'avantage sur les définitions classiques (voir par exemple [2][17]), de ne pas nécessiter de déformation du complexe elliptique considéré.

Les §§1 et 2 consistent en rappels sur la notion classique de «~structure mixte~»\footnote{On conseille au lecteur averti de les sauter.}. Nous ne donnons, bien entendu, que le minimum de définitions et résultats requis pour pouvoir énoncer et démontrer, au §3, le théorème central que nous avons en vue, et qui dit, de manière un peu plus précise, que l'image directe, par la projection (supposée propre) sur la base (supposée localement compacte), d'un complexe elliptique relatif est un complexe parfait. Quand la base $S$ est compacte, la définition de l'indice analytique en tant qu'élément de $K(S)$ («~groupe de Grothendieck~» des fibrés vectoriels complexes de rang fini sur $S$) découle alors du sorite de globalisation de (II 2.3.2).

Dans toute la suite, si $S$ est un espace topologique, on notera $\cO_S$ le faisceau des fonctions continues sur $S$ à valeurs réelles.

\subsection*{1. Le sorite des structures mixtes}

Dans ce numéro, on se fixe un espace topologique $S$.

\subsubsection*{1.1. Définition des variétés mixtes}
Soit $V$ un ouvert de $\R^n$. Pour $r\in\N$, on note $C^r(V)$ l'espace de Fréchet des fonctions de classe $C^r$ sur $V$ à valeurs réelles. On pose
\[
  C^\infty(V)=\varprojlim_r C^r(V).
\]
Soit $U$ un ouvert de $S$. Pour $r\in\N\cup\{\infty\}$, on note ${}^r\cO_{U\times V}$ le faisceau d'anneaux sur $U\times V$ défini par
